\section{Related Work}
\subsection{LLM-Based Multi-Agent Collaboration}
Large language models (LLMs) have recently demonstrated impressive capabilities as versatile task-solving agents, attracting substantial interest in both industry and academia \cite{wang2023survey,xi2023rise}. Since LLMs are trained on natural language corpus that is biased toward human thinking \cite{zhuge2023mindstorms} and optimized for conversation \cite{ouyang2022training}, LLM-based agents can collaborate through natural language in a human-like manner and harness a range of benefits that come with collaboration.

Recent works have initiated attempts to coordinate agents with varied expertise in order to improve outcomes on a wide spectrum of tasks that benefit from a diversity of knowledge. Medagent \cite{tang2023MedAgents} collects medical agents with different specialties to provide a comprehensive analysis of patient's conditions and treatment options. MetaGPT \cite{hong2023metagpt} and ChatDev \cite{chetDev} enable agents with different roles such as product managers, designers, and programmers to collaborate in software development, thereby improving the quality of the software produced. MARG\cite{darcy2024marg} develops a framework for integrating the proficiency of multiple expert agents to review scientific papers. AutoAgents\cite{AutoAgents} and OKR-Agent\cite{OKRagents} demonstrate how creative content tasks, such as creative writing and storyboard generation, can benefit from the collaboration of agents with diverse domain backgrounds. 
AgentVerse\cite{chen2023agentverse} showcases a scenario in which multiple experts with different backgrounds collaborate to deliver a hydrogen storage station siting solution. 

Additionally, recent studies have discovered that multiple agents can work together to foster cognitive synergy \cite{luppi2022synergistic} similar to humans. Chan et al. request multiple agents to delve into the discussion from diverse perspectives to catalyze a comprehensive assessment that is greater than the sum of their separate assessments. Liang et al.\cite{Encouraging_Divergent_Thinking_Debate} and Du et al.\cite{du2023improving_debate} let multiple agents to debate with each other to encourage deeper levels of contemplation. Zhuge et al.\cite{zhuge2023mindstorms} propose the concept of ``mindstorm'' to describe how multiple agents take multiple rounds of communication to iterate the ideas to find a solution that is often superior to any individual solution. 

Despite the potential of multi-agent collaboration shown in various fields, most works require coding to design coordination strategies for agents, limiting accessibility for broader general users. While AutoGen \cite{wu2023autogen} and AutoAgents \cite{AutoAgents} support representing coordination strategies in pure natural language, users still encounter a series of issues when using natural language to explore and design coordination strategies. Our work attempts to address these issues with structured generation of coordination strategies and a set of visualization approaches to facilitate users' understanding and exploration of the coordination strategies.


\subsection{Generating Coordination Strategy using LLMs}
Although LLM-based agents have shown tremendous potential in collaboratively completing tasks, manually designing coordination strategies is often challenging, time-consuming, and sometimes requires expertise in specific domain knowledge \cite{li2023camel}. Thus, it is highly desirable to leverage LLM's inherent coordinating capabilities and prior knowledge across different tasks to aid in designing coordination strategies for agent collaboration.

Many works on multi-agent collaboration leverage the prior knowledge of LLMs for the formation and adjustment of agent teams. Wang et al. \cite{wang2023unleashing} prompt LLM to dynamically identify a set of agent roles in response to a task query. Medagent \cite{tang2023MedAgents} prompts LLM to work as a medical expert who specializes in categorizing a specific medical scenario into specific
areas of medicine and gathering corresponding expert agents. DyLAN \cite{liu2023dynamic} leverages LLM to score the performance of agents and dynamically optimize the team organization during collaboration. The AgentBuilder module of Autogen \cite{wu2023autogen} prompts LLM to generate system prompts for multiple agents based on the current task and add them to a group chat for collaboration. AutoAgents \cite{AutoAgents} designs an LLM-based Agent Observer to check the compliance of the agent with the requirements and make suggestions for adjustments.

LLMs are also widely utilized to aid in planning the collaboration process of multiple agents. For example, AutoAgents \cite{AutoAgents} prompt LLMs to draft a collaboration plan that specifies the agents involved in each step and the expected outputs. In Autogen's group chat mode \cite{wu2023autogen}, the user can play the role of an Admin agent to draft and refine the collaboration plan together with an LLM-based Planner agent. OKR-agent \cite{OKRagents} leverage LLMs to recursively decompose the tasks for teams of agents. Further, AgentVerse \cite{chen2023agentverse} designs a collaborative decision-making stage for multiple LLM-based agents to make short-term planning.

Our work focuses more on how to leverage LLM to facilitate general users in designing their own multi-agent coordination strategy. To achieve this, we propose an LLM-based three-stage generation method to generate a structured coordination strategy based on the user's goal. Additionally, we propose a set of interactions to assist users in flexibly exploiting the coordination ability of LLMs during exploration.


\subsection{Interface for LLM-based Agents}
During the execution of LLM-based agents, a multitude of intricate information is involved, which is hard to digest with a plain text terminal \cite{weng2024insightlens,lu2024agentlens}. Therefore, it is highly desirable to have some interfaces to assist in understanding and intervening in the execution process. Early interfaces \cite{AgentGPT, xagent2023, SuperAGI} for monitoring single LLM-based agents typically feature an outline view that maps out the overall execution process, complemented by detailed text blocks enhanced with highlighting and icons. For systems\cite{park2023generative, lin2023agentsims, CooperativeEmbodiedAgent, chetDev} that deploys multiple agents in a virtual sandbox environment, a panoramic view is usually provided to transform text-form information into concrete visual elements (e.g. moving agent avatars, expressive emojis) for easier comprehension of the overall process. Topological structures such as trees and graphs are also utilized to visualize and manage the execution processes of agents. SPROUT\cite{liu2023sprout} employs a tree structure to assist users in visualizing and controlling the process of an agent composing code tutorials. Hong et al. \cite{hong2024data} use a hierarchical graph structure to manage the execution process of a data science agent and allow users to interactively edit the graph during execution. AutoGen \cite{wu2023autogen} introduces a transition graph to allow users to constrain agent transition to mitigate the risk of sub-optimal agent transitions during multi-agent collaboration in Group Chat mode. Recently, AgentLens \cite{lu2024agentlens} initiated the first attempt to design a visual analysis system to assist users in analyzing the agent behaviors in LLM-based multi-agent systems.

Extending this line of work, our work enables general users to visually explore coordination strategies for LLM-based multi-agent collaboration. 



